\chapter{Marco teórico}
\label{cap:marco_teorico}

El Método de los Elementos Finitos (MEF) es el procedimiento numérico de referencia para aproximar la solución de problemas de mecánica de medios continuos formulados como ecuaciones diferenciales en derivadas parciales. Su solidez teórica está documentada en los textos clásicos de la disciplina \autocite{zienkiewicz2013fem,bathe2014fem,cook2002concepts}. Este capítulo abre situando el trabajo en su contexto —los antecedentes y el estado del arte del software para la enseñanza del MEF, que justifican la oportunidad de EduFEM— y continúa estableciendo los fundamentos matemáticos que sustentan su motor de cálculo, restringido a problemas de elasticidad lineal plana resueltos con elementos isoparamétricos cuadriláteros de cuatro y nueve nodos. La exposición de esos fundamentos sigue el orden del propio canal de cálculo del programa: de la formulación variacional a la recuperación de tensiones, pasando por la calidad de malla y la verificación numérica.

\section{Antecedentes y estado del arte}
\label{sec:antecedentes}

El MEF se enseña tradicionalmente combinando el desarrollo analítico en pizarra con el uso de paquetes comerciales de propósito general. Esa práctica plantea una tensión pedagógica persistente: las herramientas profesionales son cajas negras optimizadas para producir resultados, no para revelar el procedimiento intermedio, mientras que el desarrollo puramente manual no escala más allá de uno o dos elementos. La investigación educativa documenta esta brecha y propone el software específicamente didáctico como puente entre ambos extremos \autocite{bishay2020teaching,lee2015interactive,perezsantiago2023fem}. Un estudio que recoge la opinión de expertos de la industria y de la academia sobre la enseñanza del método señala, además, que el diseño de estos cursos tiende a privilegiar las destrezas prácticas por encima de los conceptos fundamentales, y que el egresado no solo debe saber operar un modelo, sino también planificarlo, verificarlo y validarlo \autocite{perezsantiago2023fem}.

Existen diversos antecedentes de software con vocación didáctica, que pueden agruparse en tres familias. La primera reúne las herramientas centradas en el análisis matricial de estructuras de barras —pórticos y reticulados planos—: Bishay propone un constructor y un analizador de reticulados con los que el estudiante arma estructuras de barras y observa el ensamblaje y la resolución del sistema, complementados con miniproyectos generados por los propios alumnos \autocite{bishay2020teaching}, en la línea de programas previos como ED-Beams y ED-Frames \autocite{suarez1998edelas2d}; estas herramientas, sin embargo, no abordan el continuo elástico bidimensional. La segunda familia la integran los simuladores que profundizan en un eslabón aislado del método: Lee desarrolla simuladores interactivos del procesamiento de las ecuaciones de elementos finitos \autocite{lee2015interactive} y de la visualización de los modos propios de la matriz de rigidez —modos de sólido rígido, modos espurios de energía nula (\emph{hourglass}) y bloqueo por cortante, no modos de vibración— \autocite{lee2015eigenmodes}. Ambos son funciones educativas embebidas en VisualFEA, un programa comercial de análisis por elementos finitos. La tercera reúne las herramientas de cálculo de propósito general, que no fueron concebidas como material didáctico y que dejan el procedimiento fuera del alcance visual del estudiante, aunque por motivos opuestos. Por un lado, las bibliotecas de programación —CALFEM o FEniCS, entre otras—: en ellas el cálculo es enteramente inspeccionable, pero en el código fuente y no en una representación gráfica, de modo que exigen una competencia de programación previa y no ofrecen visualización guiada sobre el modelo. Por otro, los paquetes comerciales de análisis estructural \autocite{csi2017sap2000}, que operan desde la perspectiva del estudiante como cajas negras: entregan resultados sin exponer las matrices intermedias que los produjeron. Para estas dos categorías, la comparación de la \autoref{tab:comparativa} se restringe deliberadamente a atributos que se desprenden de su propia naturaleza —presencia o ausencia de interfaz gráfica, régimen de licencia, requisitos de ejecución y exposición del cálculo intermedio— y no a juicios sobre la calidad didáctica de cada producto; los atributos que exigirían examinar cada producto en particular se dejan sin evaluar y se marcan con una raya. VisualFEA es, de hecho, el producto comercial con orientación didáctica declarada más mencionado en la literatura, y es el programa en el que están embebidos los simuladores de Lee ya descritos; al ser cerrado y de pago no se lo pudo examinar de primera mano, de modo que lo consignado aquí procede de las publicaciones de sus propios autores y se incluye para no presentar el vacío como más amplio de lo que esta revisión permite sostener.

El antecedente más directo de este trabajo es ED-Elas2D, un programa educativo desarrollado en la Universitat Politècnica de Catalunya para el adiestramiento asistido por computadora en el análisis de sólidos bidimensionales por el MEF \autocite{suarez1998edelas2d}. Desarrollado en una institución española y con una interfaz al menos parcialmente en castellano —a juzgar por las capturas de pantalla de su publicación, que no declara el idioma—, ED-Elas2D comparte con EduFEM no solo el dominio —elasticidad plana en tensión y deformación plana— y la organización en pre-proceso, proceso y post-proceso, sino también, en buena medida, el idioma e, incluso, una biblioteca de elementos triangulares y cuadriláteros que incluye los elementos Q4, Q8 y Q9. Su existencia confirma la pertinencia de un entorno educativo de esta clase y, a la vez, ayuda a delimitar el aporte del presente trabajo, que no reside en la lengua sino en la concepción y la tecnología. Frente a aquel programa de finales de la década de 1990 —software académico cuya licencia y disponibilidad la publicación no documenta, ligado a la plataforma y a las posibilidades de su época—, EduFEM se distribuye como software libre y de código abierto, expone de forma interactiva cada matriz del cálculo ($\bm{N}$, $\bm{J}$, $\bm{B}$, $\bm{D}$ y $\bm{k}_e$) sobre el elemento que el alumno selecciona en el lienzo, genera de manera automática una memoria de cálculo con sustitución numérica paso a paso, trazable al modelo del usuario, e incorpora una batería de verificación y validación documentada.

La \autoref{tab:comparativa} sitúa a EduFEM frente a estas categorías de recursos según los atributos relevantes para la enseñanza del continuo elástico bidimensional: los que describen el recurso —idioma, dimensión, elementos, licencia y requisitos de ejecución— y los que distinguen a este trabajo —la exposición del cálculo paso a paso, la cobertura del canal completo, la memoria de cálculo automática, la posibilidad de observar los fenómenos numéricos que la introducción señala y la verificación y validación publicadas junto con la herramienta—. Cuando las publicaciones consultadas no documentan un atributo, la celda dice «no consta»: registra un límite de la revisión, no una carencia del recurso. El contraste evidencia un vacío: ninguno reúne, de forma simultánea, una interfaz íntegramente en español, el carácter libre y de código abierto, el dominio del continuo elástico plano, la exposición interactiva del cálculo paso a paso sobre el modelo concreto del usuario, la cobertura del canal completo del MEF —del mapeo isoparamétrico a la recuperación de tensiones—, una memoria de cálculo trazable a ese modelo, la observación guiada de la distorsión de malla, la sub-integración, el bloqueo por cortante y el condicionamiento, y una verificación y validación documentadas. Solo los simuladores de VisualFEA comparten uno de estos diferenciadores: hacen visibles los modos espurios y el bloqueo a través de los modos propios de la matriz de rigidez, aunque en un producto comercial y en inglés.

\begin{table}[htbp]
\centering
\caption[Comparación cualitativa de recursos para la enseñanza del MEF frente a EduFEM.]{Comparación cualitativa de recursos para la enseñanza del MEF frente a EduFEM.}
\label{tab:comparativa}
{\footnotesize
\setlength{\tabcolsep}{2pt}
\renewcommand{\arraystretch}{1.0}
\begin{tabular}{@{}>{\raggedright\arraybackslash}p{2.0cm}>{\raggedright\arraybackslash}p{2.5cm}*{5}{>{\raggedright\arraybackslash}p{2.1cm}}@{}}
\toprule
\textbf{Atributo} & \textbf{EduFEM (este trabajo)} & \textbf{ED-Elas2D} \autocite{suarez1998edelas2d} & \textbf{Analizador de reticulados} \autocite{bishay2020teaching} & \textbf{Simuladores de VisualFEA} \autocite{lee2015interactive,lee2015eigenmodes} & \textbf{Bibliotecas de programación} (CALFEM, FEniCS) & \textbf{Software comercial} (SAP2000 \autocite{csi2017sap2000}, ANSYS, Abaqus) \\
\midrule
Idioma & Español & Español (parcial) & Inglés & Inglés & Inglés & Inglés (localización parcial) \\
Dimensión & 2D & 2D & 2D/3D (barras) & 2D/3D & 2D/3D & 2D/3D \\
Elementos & Cuadriláteros Q4 y Q9 & Triángulos y cuadriláteros (Q4/Q8/Q9) & Barras (reticulado) & Continuo, placas, láminas y sólidos & Diversos & Amplios \\
Licencia & Libre, código abierto & No consta & No declarada (código MATLAB) & Comercial & Libre & Comercial \\
Requisitos de ejecución & Instalador autocontenido (Windows) & Windows (década de 1990) & Requiere MATLAB & Requiere VisualFEA & Entorno de programación & Instalación y licencia del fabricante \\
Cálculo paso a paso & Alta: ocho módulos (M0--M7) sobre el modelo del usuario & Media (pre/proc/post) & Media-alta (código construido paso a paso) & Media-alta, en un eslabón & Visible en el código, sin visualización guiada & Baja (caja negra) \\
Canal completo & Sí & Sí & Parcial (solo barras) & Parcial (un eslabón) & Parcial (por programación) & No (no didáctico) \\
Memoria de cálculo automática & Sí: PDF con sustitución numérica, trazable al modelo del usuario & No consta & No consta & Valores intermedios en pantalla, no como documento & No (la programa el usuario) & Informes de resultados, sin el desarrollo matricial \\
Fenómenos numéricos observables & Distorsión y $\det\bm{J}$ (M0, M2); sub-integración y hourglass (M5); bloqueo por cortante Q4/Q9 (Cook); $\kappa_2$ (memoria) & No consta & No consta & Hourglass y bloqueo por cortante, vía modos propios de $\bm{K}$ & Por programación, sin visualización guiada & Sin exposición didáctica \\
V\&V publicada con la herramienta & Sí: MMS, Timoshenko y Cook (\autoref{cap:resultados}) & No consta & No consta & No consta & --- & --- \\
\bottomrule
\end{tabular}}
{\singlespacing\notaapa{«Canal completo»: el recurso recorre la cadena del MEF de extremo a extremo (mapeo isoparamétrico, Jacobiano, matriz $\bm{B}$, constitutiva, rigidez, ensamblaje y recuperación de tensiones). «No consta»: las publicaciones consultadas no documentan el atributo, lo que no implica que el recurso carezca de él. La raya marca atributos no evaluados para las categorías de propósito general, según el criterio enunciado en el texto.}}
\end{table}

En síntesis, salvo el antecedente directo de ED-Elas2D, la mayoría de estos recursos se concentra en un eslabón aislado del método —el ensamblaje, la respuesta dinámica, un tipo particular de elemento— o se limita a estructuras de barras y reticulados \autocite{bishay2020teaching}, dejando sin cubrir, con el detalle interno que requiere la enseñanza, el continuo elástico bidimensional; y ni siquiera ED-Elas2D —software académico condicionado por la tecnología de su época— reúne el conjunto de atributos que persigue este trabajo. La brecha que lo motiva es, por tanto, la ausencia de un entorno integrado, libre y de código abierto, que recorra el canal de cálculo completo del MEF en elasticidad plana —del mapeo isoparamétrico a la recuperación de tensiones— exponiendo de forma interactiva, y en español, cada matriz y cada paso numérico sobre el modelo concreto del alumno, haciendo observables los fenómenos numéricos del método, generando una memoria de cálculo trazable a ese modelo e incorporando verificación y validación rigurosas. EduFEM se diseña para cubrir precisamente ese vacío, articulando el rigor numérico que se desarrolla en el resto de este capítulo con una capa pedagógica que hace observable el procedimiento que los textos clásicos formulan en abstracto.

\section{Fundamentos del MEF y formulación variacional en elasticidad lineal}
\label{sec:formulacion-debil}

Sea un dominio bidimensional $\Omega \subset \mathbb{R}^2$ con frontera $\Gamma = \Gamma_u \cup \Gamma_t$, donde $\Gamma_u$ es la porción con desplazamientos prescritos (restricciones de Dirichlet) y $\Gamma_t$ la porción con tracciones impuestas (Neumann). El problema de equilibrio en elasticidad lineal estática se enuncia, en su forma fuerte, como

\begin{equation}
\nabla \cdot \bm{\sigma} + \bm{b} = \bm{0} \quad \text{en } \Omega,
\label{eq:equilibrio-fuerte}
\end{equation}

donde $\bm{\sigma}$ es el tensor de tensiones y $\bm{b}$ la fuerza por unidad de volumen. La forma fuerte exige continuidad de segundo orden del campo de desplazamientos $\bm{u}$, requisito difícil de satisfacer de manera global. El MEF parte por ello de la \emph{forma débil} o variacional, obtenida multiplicando la \autoref{eq:equilibrio-fuerte} por un campo de desplazamientos virtuales $\delta\bm{u}$ admisible (nulo sobre $\Gamma_u$) e integrando por partes mediante el teorema de la divergencia. Esto conduce al principio de los trabajos virtuales \autocite{zienkiewicz2013fem,hughes2000fem,reddy2006introduction}:

\begin{equation}
\int_\Omega \delta\bm{\varepsilon}^{T} \bm{\sigma} \, d\Omega
= \int_\Omega \delta\bm{u}^{T} \bm{b} \, d\Omega
+ \int_{\Gamma_t} \delta\bm{u}^{T} \bm{t} \, d\Gamma,
\label{eq:trabajos-virtuales}
\end{equation}

con $\delta\bm{\varepsilon}$ las deformaciones virtuales y $\bm{t}$ las tracciones de borde. La formulación débil rebaja el requisito de continuidad a primer orden, lo que habilita aproximaciones polinómicas a trozos. El campo continuo se discretiza interpolando los valores nodales mediante funciones de forma: dentro de cada elemento, $\bm{u} \approx \bm{N}\,\bm{u}_e$, donde $\bm{N}$ agrupa las funciones de forma y $\bm{u}_e$ los desplazamientos de los nodos del elemento. Sustituyendo esta aproximación en la \autoref{eq:trabajos-virtuales} y aprovechando la arbitrariedad de $\delta\bm{u}$ se obtiene el sistema algebraico de equilibrio

\begin{equation}
\bm{K}\,\bm{u} = \bm{F},
\label{eq:sistema-global}
\end{equation}

donde $\bm{K}$ es la matriz de rigidez global, $\bm{u}$ el vector de desplazamientos nodales y $\bm{F}$ el vector de cargas. El resto del capítulo desarrolla cada pieza de la \autoref{eq:sistema-global}.

\section{Elasticidad plana: tensión plana y deformación plana}
\label{sec:elasticidad-plana}

La elasticidad bidimensional admite dos idealizaciones de un sólido tridimensional \autocite{timoshenko1970elasticity,cook2002concepts}. La \emph{tensión plana} modela cuerpos delgados cargados en su plano, en los que la tensión normal al plano es despreciable ($\sigma_z = 0$); es el caso de una chapa o membrana de poco espesor. La \emph{deformación plana} modela cuerpos prismáticos largos en los que la deformación longitudinal se anula ($\varepsilon_z = 0$); corresponde, por ejemplo, a una presa o un tubo largo donde una rebanada representa todo el sólido.

En ambos casos las componentes relevantes se agrupan en notación de Voigt como $\bm{\sigma} = [\sigma_x,\, \sigma_y,\, \tau_{xy}]^T$ y $\bm{\varepsilon} = [\varepsilon_x,\, \varepsilon_y,\, \gamma_{xy}]^T$, ligadas por la ley constitutiva $\bm{\sigma} = \bm{D}\,\bm{\varepsilon}$. Para tensión plana la matriz constitutiva es

\begin{equation}
\bm{D}_{\text{TP}} = \frac{E}{1-\nu^{2}}
\begin{bmatrix}
1 & \nu & 0 \\
\nu & 1 & 0 \\
0 & 0 & \dfrac{1-\nu}{2}
\end{bmatrix},
\label{eq:D-tension-plana}
\end{equation}

y para deformación plana

\begin{equation}
\bm{D}_{\text{DP}} = \frac{E}{(1+\nu)(1-2\nu)}
\begin{bmatrix}
1-\nu & \nu & 0 \\
\nu & 1-\nu & 0 \\
0 & 0 & \dfrac{1-2\nu}{2}
\end{bmatrix},
\label{eq:D-deformacion-plana}
\end{equation}

donde $E$ es el módulo de elasticidad y $\nu$ el coeficiente de Poisson. El cambio en el denominador refleja el acoplamiento adicional que introduce la restricción $\varepsilon_z = 0$. La deformación plana es sensible a la incompresibilidad: cuando $\nu \to 0{,}5$ el factor $(1-2\nu)$ tiende a cero y $\bm{D}_{\text{DP}}$ se degrada, fenómeno asociado al bloqueo volumétrico (\emph{volumetric locking}) que se analiza en la \autoref{sec:locking}.

\section{Elementos isoparamétricos Q4 y Q9}
\label{sec:isoparametricos}

EduFEM emplea elementos isoparamétricos cuadriláteros, en los que la misma familia de funciones de forma interpola tanto la geometría como el campo de desplazamientos \autocite{bathe2014fem,onate2009structural}. Las funciones se definen sobre un elemento de referencia (el cuadrado natural $(\xi, \eta) \in [-1, 1]^2$), y un mapeo isoparamétrico las lleva a la geometría física.

El elemento \textbf{Q4} es un cuadrilátero de cuatro nodos con interpolación bilineal. Numerados sus vértices en sentido antihorario, las funciones de forma son

\begin{equation}
\begin{aligned}
N_1 &= \tfrac{1}{4}(1-\xi)(1-\eta), & N_2 &= \tfrac{1}{4}(1+\xi)(1-\eta), \\
N_3 &= \tfrac{1}{4}(1+\xi)(1+\eta), & N_4 &= \tfrac{1}{4}(1-\xi)(1+\eta).
\end{aligned}
\label{eq:N-q4}
\end{equation}

El elemento \textbf{Q9} añade cuatro nodos medios de arista y un nodo central, alcanzando interpolación bicuadrática como producto tensorial de polinomios de Lagrange unidimensionales. Sus funciones de forma para vértices, lados y centro son, respectivamente,

\begin{equation}
\begin{aligned}
N_1 &= \tfrac{1}{4}\,\xi(\xi-1)\,\eta(\eta-1), & N_2 &= \tfrac{1}{4}\,\xi(\xi+1)\,\eta(\eta-1), \\
N_3 &= \tfrac{1}{4}\,\xi(\xi+1)\,\eta(\eta+1), & N_4 &= \tfrac{1}{4}\,\xi(\xi-1)\,\eta(\eta+1), \\
N_5 &= \tfrac{1}{2}(1-\xi^2)\,\eta(\eta-1), & N_6 &= \tfrac{1}{2}\,\xi(\xi+1)(1-\eta^2), \\
N_7 &= \tfrac{1}{2}(1-\xi^2)\,\eta(\eta+1), & N_8 &= \tfrac{1}{2}\,\xi(\xi-1)(1-\eta^2), \\
N_9 &= (1-\xi^2)(1-\eta^2). &&
\end{aligned}
\label{eq:N-q9}
\end{equation}

Todas satisfacen la propiedad de Kronecker $N_i(\xi_j, \eta_j) = \delta_{ij}$ y forman una partición de la unidad. El mapeo isoparamétrico interpola la posición física de cualquier punto interior a partir de las coordenadas nodales $(x_i, y_i)$:

\begin{equation}
x(\xi, \eta) = \sum_{i=1}^{n} N_i(\xi, \eta)\,x_i, \qquad
y(\xi, \eta) = \sum_{i=1}^{n} N_i(\xi, \eta)\,y_i,
\label{eq:mapeo-iso}
\end{equation}

con $n = 4$ para Q4 y $n = 9$ para Q9. El elemento Q9, de orden superior, captura gradientes de tensión más complejos y resiste mejor las distorsiones geométricas, a costa de un mayor número de grados de libertad (GDL): 8 GDL en Q4 frente a 18 GDL en Q9. La \autoref{fig:isoparametricos} ilustra ambos elementos en el cuadrado natural y su mapeo al espacio físico.

\begin{figure}[htbp]
\centering
\includegraphics[width=\linewidth]{fig_isoparametricos}
\caption{Elementos isoparamétricos Q4 y Q9 y el mapeo isoparamétrico al espacio físico.}
\label{fig:isoparametricos}
\end{figure}

\section{Jacobiano del mapeo}
\label{sec:jacobiano}

El mapeo de la \autoref{eq:mapeo-iso} relaciona los dos espacios mediante la matriz Jacobiana, que transforma derivadas respecto de las coordenadas naturales en derivadas respecto de las físicas:

\begin{equation}
\bm{J} =
\begin{bmatrix}
\dfrac{\partial x}{\partial \xi} & \dfrac{\partial y}{\partial \xi} \\[6pt]
\dfrac{\partial x}{\partial \eta} & \dfrac{\partial y}{\partial \eta}
\end{bmatrix}
= \frac{\partial \bm{N}}{\partial(\xi,\eta)}\,\bm{X}_e,
\label{eq:jacobiano}
\end{equation}

donde $\partial\bm{N}/\partial(\xi,\eta)$ es la matriz $2\times n$ cuyas filas son las derivadas $\partial N_i/\partial\xi$ y $\partial N_i/\partial\eta$, y $\bm{X}_e$ ($n\times 2$) reúne las coordenadas nodales del elemento. Con esta convención —la de Bathe \autocite[p.~345]{bathe2014fem}, y la que implementa el motor— cada fila de $\bm{J}$ es el vector tangente a una línea de coordenada natural ($\partial\bm{x}/\partial\xi$ y $\partial\bm{x}/\partial\eta$), y la regla de la cadena de la \autoref{eq:fisica-natural} emplea $\bm{J}^{-1}$ directamente; la convención transpuesta, también frecuente en la literatura, exigiría $\bm{J}^{-\mathsf{T}}$. Su determinante mide el factor local de cambio de área entre ambos espacios, $dx\,dy = |\det\bm{J}|\,d\xi\,d\eta$, y por tanto es central en la integración numérica. Para una matriz $2\times2$ se calcula de forma explícita,

\begin{equation}
\det\bm{J} = J_{11}J_{22} - J_{12}J_{21},
\label{eq:detJ}
\end{equation}

y su inversa se obtiene por cofactores, evitando rutinas genéricas de inversión:

\begin{equation}
\bm{J}^{-1} = \frac{1}{\det\bm{J}}
\begin{bmatrix}
J_{22} & -J_{12} \\
-J_{21} & J_{11}
\end{bmatrix}.
\label{eq:invJ}
\end{equation}

La condición $\det\bm{J} > 0$ en todo el elemento garantiza que el mapeo sea biyectivo y el elemento válido \autocite{bathe2014fem,cook2002concepts}. Un determinante nulo o negativo indica un elemento degenerado o invertido; EduFEM impone un umbral mínimo ($\texttt{JACOBIAN\_MIN\_DETERMINANT}$) para detectar y rechazar esas geometrías antes de resolver.

\section{Matriz de deformación-desplazamiento B}
\label{sec:matriz-b}

Las deformaciones se obtienen derivando el campo de desplazamientos interpolado. La relación cinemática en elasticidad plana es

\begin{equation}
\bm{\varepsilon} =
\begin{bmatrix} \varepsilon_x \\ \varepsilon_y \\ \gamma_{xy} \end{bmatrix}
=
\begin{bmatrix}
\dfrac{\partial u}{\partial x} \\[4pt]
\dfrac{\partial v}{\partial y} \\[4pt]
\dfrac{\partial u}{\partial y} + \dfrac{\partial v}{\partial x}
\end{bmatrix},
\label{eq:cinematica}
\end{equation}

que en forma matricial se escribe $\bm{\varepsilon} = \bm{B}\,\bm{u}_e$. La matriz $\bm{B}$, de dimensión $3 \times 2n$, contiene las derivadas físicas de las funciones de forma:

\begin{equation}
\bm{B} =
\begin{bmatrix}
\dfrac{\partial N_1}{\partial x} & 0 & \dfrac{\partial N_2}{\partial x} & 0 & \cdots \\[6pt]
0 & \dfrac{\partial N_1}{\partial y} & 0 & \dfrac{\partial N_2}{\partial y} & \cdots \\[6pt]
\dfrac{\partial N_1}{\partial y} & \dfrac{\partial N_1}{\partial x} & \dfrac{\partial N_2}{\partial y} & \dfrac{\partial N_2}{\partial x} & \cdots
\end{bmatrix}.
\label{eq:matriz-B}
\end{equation}

Las funciones de forma se definen en coordenadas naturales, pero $\bm{B}$ exige sus derivadas físicas $\partial N_i/\partial x$ y $\partial N_i/\partial y$. El puente lo provee la inversa del Jacobiano, aplicada por columna a cada función de forma según la regla de la cadena:

\begin{equation}
\begin{bmatrix}
\dfrac{\partial N_i}{\partial x} \\[6pt]
\dfrac{\partial N_i}{\partial y}
\end{bmatrix}
= \bm{J}^{-1}
\begin{bmatrix}
\dfrac{\partial N_i}{\partial \xi} \\[6pt]
\dfrac{\partial N_i}{\partial \eta}
\end{bmatrix}.
\label{eq:fisica-natural}
\end{equation}

Esta relación vectorial es la conexión esencial entre el espacio natural, donde las derivadas se conocen analíticamente, y el espacio físico, donde se formula la deformación.

\section{Matriz de rigidez elemental e integración numérica de Gauss}
\label{sec:rigidez-gauss}

La matriz de rigidez de un elemento integra el producto $\bm{B}^T\bm{D}\bm{B}$ sobre su dominio físico, con $t$ el espesor:

\begin{equation}
\bm{k}_e = \int_{\Omega_e} \bm{B}^T \bm{D}\, \bm{B}\, t \; dA
= \int_{-1}^{1}\!\!\int_{-1}^{1} \bm{B}^T(\xi,\eta)\,\bm{D}\,\bm{B}(\xi,\eta)\, t\, |\det\bm{J}|\; d\xi\, d\eta.
\label{eq:ke-integral}
\end{equation}

El integrando no es polinómico: como $\bm{B}$ depende de $\bm{J}^{-1}$, que es racional en $(\xi,\eta)$, la integración algebraica exacta deja de ser practicable salvo en los elementos más simples \autocite{zienkiewicz2013fem}. Por ello se aproxima mediante cuadratura de Gauss-Legendre \autocite{hughes2000fem}, que evalúa el integrando en puntos óptimos y lo pondera:

\begin{equation}
\bm{k}_e \approx \sum_{p=1}^{n_g} w_p\; \bm{B}^T(\xi_p,\eta_p)\,\bm{D}\,\bm{B}(\xi_p,\eta_p)\, t\, |\det\bm{J}(\xi_p,\eta_p)|,
\label{eq:ke-gauss}
\end{equation}

donde los puntos $(\xi_p,\eta_p)$ y pesos $w_p$ surgen del producto tensorial de la regla unidimensional. La regla de $m$ puntos integra exactamente polinomios de grado $2m-1$ \autocites[pp.~141-142]{hughes2000fem}[p.~161]{zienkiewicz2013fem}. EduFEM usa cuadratura $2\times2$ para Q4 (cuatro puntos en $\xi = \pm 1/\sqrt{3}$) y $3\times3$ para Q9 (puntos en $\xi = \pm\sqrt{3/5},\,0$ con pesos $5/9,\,8/9,\,5/9$), órdenes que integran exactamente la rigidez de los respectivos elementos no distorsionados.

El orden de integración no es arbitrario. Una \emph{subintegración} puede dejar la matriz de rigidez con rango deficiente y admitir modos de deformación de energía nula: los modos espurios (\emph{hourglass}), patrones de deformación no físicos que el elemento no puede resistir \autocite{cook2002concepts,bathe2014fem}. Una \emph{sobreintegración} (por ejemplo $3\times3$ en Q4) no aporta precisión apreciable —el integrando solo es polinómico en los elementos no distorsionados, y en los distorsionados la diferencia es marginal— y solo encarece el cálculo. La elección $2\times2$/$3\times3$ equilibra exactitud y estabilidad para esta familia de elementos.

\section{Fenómenos numéricos en elementos de bajo orden: bloqueo y modos espurios}
\label{sec:locking}

La discretización con elementos de bajo orden introduce fenómenos numéricos que no provienen de la física del problema, sino de las limitaciones de la interpolación, y cuya comprensión resulta esencial para interpretar correctamente los resultados \autocite{cook2002concepts,hughes2000fem,bathe2014fem}. Tres son particularmente relevantes para los cuadriláteros que emplea EduFEM y reaparecen, ya cuantificados, en el \autoref{cap:resultados}.

El \textbf{bloqueo por cortante} (\emph{shear-locking}) afecta al elemento bilineal Q4 en estados dominados por flexión. Bajo flexión pura, una sección debería curvarse sin desarrollar deformación de corte; sin embargo, el campo de desplazamientos bilineal del Q4 no puede representar esa curvatura sin generar al mismo tiempo deformaciones de corte espurias, ausentes en la solución exacta. Esas deformaciones parásitas absorben energía y vuelven al elemento artificialmente rígido, de modo que subestima los desplazamientos y converge con lentitud al refinar la malla. El elemento Q9, cuya interpolación bicuadrática sí reproduce el modo de curvatura, está prácticamente libre de este bloqueo \autocite{cook2002concepts,hughes2000fem}. La membrana de Cook (\autoref{sec:verificacion-validacion}) es el banco de prueba clásico para evidenciarlo.

El \textbf{bloqueo volumétrico} (\emph{volumetric locking}) aparece en deformación plana cuando el material se aproxima a la incompresibilidad ($\nu \to 0{,}5$). Como anticipó la \autoref{sec:elasticidad-plana}, el factor $(1-2\nu)$ del denominador de $\bm{D}_{\text{DP}}$ tiende a cero, y la restricción de incompresibilidad sobre-condiciona el elemento: la malla no dispone de grados de libertad suficientes para satisfacer la condición $\varepsilon_x + \varepsilon_y \to 0$ en cada punto sin rigidizarse en exceso. El resultado es, de nuevo, una respuesta artificialmente rígida, acompañada a menudo de oscilaciones espurias en el campo de tensiones \autocite{hughes2000fem,bathe2014fem}.

Los \textbf{modos espurios} (\emph{hourglass}) son el fenómeno inverso. Como se anticipó en la \autoref{sec:rigidez-gauss}, surgen al \emph{subintegrar} la rigidez —por ejemplo, evaluando el Q4 con un único punto de Gauss—: la matriz elemental queda con rango deficiente y admite modos de deformación de energía nula, patrones no físicos que contaminan la solución \autocite{cook2002concepts,bathe2014fem}. EduFEM los evita empleando la integración completa de cada elemento ($2\times2$ para Q4 y $3\times3$ para Q9).

Existen técnicas correctivas para los dos primeros bloqueos —entre ellas la integración reducida selectiva (SRI), que subintegra solo los términos responsables del bloqueo, y la formulación $\bar{\bm{B}}$ (\emph{B-bar}), que proyecta la parte volumétrica o de corte de la deformación \autocite{hughes2000fem,bathe2014fem}—. EduFEM no las implementa: ofrece en su lugar el elemento Q9, que evita el bloqueo por cortante por la vía del mayor orden de interpolación —no así el volumétrico, frente al cual el bicuadrático con integración completa también requiere integración selectiva \autocite[p.~222]{hughes2000fem}—, y aprovecha el comportamiento del Q4 como recurso didáctico, para que el estudiante observe el fenómeno en lugar de ocultárselo. Su incorporación se plantea como línea de extensión futura entre las recomendaciones del trabajo.

\section{Ensamblaje, restricciones y cargas}
\label{sec:ensamblaje}

Las matrices y vectores elementales se ensamblan en el sistema global mediante el mapeo de cada GDL local a su índice global. La matriz de rigidez global se obtiene sumando las contribuciones,

\begin{equation}
\bm{K} = \Ensamblaje_{e} \bm{k}_e, \qquad
\bm{F} = \sum_e \bm{f}_e^{\text{vol}} + \sum_n \bm{F}_n + \sum_\Gamma \bm{F}_\Gamma,
\label{eq:ensamblaje}
\end{equation}

donde un nodo compartido por varios elementos acumula los aportes de todos ellos. Para mallas grandes $\bm{K}$ es dispersa (mayoritariamente nula), por lo que EduFEM la construye en formato de coordenadas (COO) y la convierte a filas comprimidas (CSR) para el solucionador, apoyándose en el módulo de matrices dispersas de SciPy \autocite[p.~265]{virtanen2020scipy}.

El vector de cargas reúne tres aportes. Las \textbf{cargas nodales} $\bm{F}_n$ se aplican directamente en los GDL correspondientes. Las \textbf{cargas volumétricas} $\bm{f}_e^{\text{vol}}$ —como la gravedad— se distribuyen consistentemente integrando $\int_{\Omega_e} \bm{N}^T \bm{b}\, t\, dA$ por cuadratura de Gauss. Las \textbf{cargas superficiales} $\bm{F}_\Gamma$ sobre una arista se convierten en fuerzas nodales equivalentes ponderando la carga distribuida con las funciones de forma. Para una arista lineal de Q4 de longitud $L$ con carga $q$ que varía linealmente entre $q_s$ y $q_e$ en sus extremos, el reparto exacto es $F_1 = \tfrac{L}{6}(2q_s + q_e)$ y $F_2 = \tfrac{L}{6}(q_s + 2q_e)$ —que con $q$ constante se reduce a $qL/2$ en cada nodo—, mientras que para la arista cuadrática de Q9 con $q$ constante la distribución resulta $\tfrac{qL}{6},\,\tfrac{4qL}{6},\,\tfrac{qL}{6}$, concentrando la mayor parte en el nodo medio \autocite[p.~113]{cook2002concepts}.

Las restricciones de Dirichlet se imponen partiendo el sistema en GDL libres ($f$) y restringidos ($r$):

\begin{equation}
\begin{bmatrix} \bm{K}_{ff} & \bm{K}_{fr} \\ \bm{K}_{rf} & \bm{K}_{rr} \end{bmatrix}
\begin{bmatrix} \bm{u}_f \\ \bm{u}_r \end{bmatrix}
=
\begin{bmatrix} \bm{F}_f \\ \bm{F}_r \end{bmatrix}.
\label{eq:particion}
\end{equation}

Para \emph{restricciones homogéneas} ($\bm{u}_r = \bm{0}$) basta resolver el bloque libre $\bm{K}_{ff}\,\bm{u}_f = \bm{F}_f$. Para \emph{restricciones no homogéneas} ($\bm{u}_r \neq \bm{0}$), necesarias en la verificación con soluciones manufacturadas, el sistema reducido incorpora el término de sustitución estática

\begin{equation}
\bm{K}_{ff}\,\bm{u}_f = \bm{F}_f - \bm{K}_{fr}\,\bm{u}_r,
\label{eq:dirichlet-no-homogeneo}
\end{equation}

que preserva la simetría del bloque libre. Un modelo bien restringido —al menos tres GDL que impidan los movimientos de cuerpo rígido— hace que $\bm{K}_{ff}$ sea definida positiva y, por tanto, no singular \autocite[pp.~726-727]{bathe2014fem}.

\section{Resolución del sistema y recuperación de tensiones}
\label{sec:resolucion-tensiones}

El sistema reducido de la \autoref{eq:dirichlet-no-homogeneo} se resuelve por factorización LU directa sobre la matriz dispersa, a través de la interfaz que SciPy expone sobre la biblioteca SuperLU \autocite[p.~264]{virtanen2020scipy}. EduFEM selecciona el ordenamiento de columnas de mínimo grado sobre $\bm{K}^{T}+\bm{K}$ (opción \texttt{MMD\_AT\_PLUS\_A} de \texttt{scipy.sparse.linalg.splu}, documentada en la propia biblioteca), que aprovecha la simetría de la matriz para reducir el llenado (\emph{fill-in}). Conocido el vector $\bm{u}$, las reacciones en los GDL restringidos se recuperan como $\bm{R} = \bm{K}_{rf}\,\bm{u}_f + \bm{K}_{rr}\,\bm{u}_r - \bm{F}_r$.

Las tensiones se calculan primero en los puntos de Gauss, evaluando $\bm{\sigma}_p = \bm{D}\,\bm{B}_p\,\bm{u}_e$. A partir de las componentes cartesianas se obtienen las tensiones principales \autocite[pp.~12-14]{timoshenko1970elasticity} y la equivalente de von Mises, la medida escalar con que habitualmente se resume un estado tensional biaxial \autocite[p.~117]{cook2002concepts}:

\begin{equation}
\begin{aligned}
\sigma_{1,2} &= \frac{\sigma_x + \sigma_y}{2} \pm \sqrt{\left(\frac{\sigma_x - \sigma_y}{2}\right)^2 + \tau_{xy}^2}, \\
\sigma_{\mathrm{VM}} &= \sqrt{\tfrac{1}{2}\left[(\sigma_1-\sigma_2)^2 + (\sigma_2-\sigma_z)^2 + (\sigma_z-\sigma_1)^2\right]}, \\
\sigma_z &=
\begin{cases}
0 & \text{tensión plana,}\\
\nu\,(\sigma_x+\sigma_y) & \text{deformación plana.}
\end{cases}
\end{aligned}
\label{eq:principales-vm}
\end{equation}

En tensión plana $\sigma_z = 0$ y la \autoref{eq:principales-vm} se reduce a la forma plana $\sigma_{\mathrm{VM}} = \sqrt{\sigma_1^2 - \sigma_1\sigma_2 + \sigma_2^2}$. En deformación plana la tensión fuera del plano no es nula y es la tercera tensión principal: omitirla —simplificación a la que invita el carácter bidimensional del planteo— subestima la tensión equivalente, por lo que EduFEM la incorpora en todas sus rutas de cálculo (puntos de Gauss, nodos extrapolados, promedio nodal, sonda y memoria de cálculo), con la $\nu$ del material de cada elemento.

En el elemento Q4 los puntos de la regla $2\times2$ coinciden con los \emph{puntos de Barlow}, ubicaciones donde las tensiones convergen con un orden superior al del resto del elemento (superconvergencia) \autocite[p.~231]{cook2002concepts}; en el Q9 los puntos óptimos son también los de la regla $2\times2$, y no los nueve de la regla $3\times3$ con que se integra su rigidez. EduFEM evalúa las tensiones en los puntos de la cuadratura de rigidez de cada elemento y las lleva a los nodos mediante una matriz precalculada, $\bm{\sigma}_{\text{nodos}} = \bm{E}\,\bm{\sigma}_{\text{Gauss}}$, en lugar de evaluar $\bm{B}$ directamente en los nodos. Para el Q4 esa elección aprovecha la superconvergencia; para el Q9 privilegia que la matriz $\bm{M}$ definida a continuación sea cuadrada con los nueve nodos del elemento, a costa de no muestrear en los puntos óptimos, y la verificación del \autoref{cap:resultados} muestra que el campo recuperado converge de todos modos a $\mathcal{O}(h^{2})$.

La matriz de extrapolación $\bm{E}$ se deduce de la propia interpolación isoparamétrica. Dentro del elemento, el campo de tensiones se aproxima con las mismas funciones de forma que la geometría, de modo que el valor en un punto de Gauss es la interpolación de los valores nodales,

\begin{equation}
\bm{\sigma}(\xi_g,\eta_g) = \sum_{i=1}^{n} N_i(\xi_g,\eta_g)\,\bm{\sigma}_{\text{nodo},i}.
\label{eq:interp-tension}
\end{equation}

Evaluando la \autoref{eq:interp-tension} en los puntos de Gauss del elemento se obtiene el sistema $\bm{\sigma}_{\text{Gauss}} = \bm{M}\,\bm{\sigma}_{\text{nodos}}$, donde $M_{ji} = N_i(\xi_j^{G},\eta_j^{G})$ reúne las funciones de forma evaluadas en cada punto de Gauss. Cuando el número de puntos de Gauss iguala al de nodos del elemento —cuadratura $2\times2$ para Q4 y $3\times3$ para Q9—, la matriz $\bm{M}$ es cuadrada e invertible, y la extrapolación es su inversa,

\begin{equation}
\bm{\sigma}_{\text{nodos}} = \bm{M}^{-1}\,\bm{\sigma}_{\text{Gauss}} = \bm{E}\,\bm{\sigma}_{\text{Gauss}}.
\label{eq:extrap-inversa}
\end{equation}

Para el elemento Q4 esta inversa admite forma cerrada. Como los cuatro puntos de Gauss se ubican en $(\pm 1/\sqrt{3},\,\pm 1/\sqrt{3})$ y los nodos en $(\pm 1,\,\pm 1)$, extrapolar equivale a evaluar la interpolación bilineal en coordenadas escaladas por $\sqrt{3}$, con entradas de la forma $\tfrac{1}{4}(1\pm\sqrt{3})(1\pm\sqrt{3})$. Con los puntos de Gauss numerados en el orden que produce el motor —$(-,-)$, $(-,+)$, $(+,-)$, $(+,+)$ en $(\xi,\eta)$— y los nodos en sentido antihorario desde $(-1,-1)$, resulta

\begin{equation}
\bm{E}_{\text{Q4}} =
\begin{bmatrix}
a & b & b & c \\
b & c & a & b \\
c & b & b & a \\
b & a & c & b
\end{bmatrix},
\qquad
a = 1+\frac{\sqrt{3}}{2},\quad
b = -\frac{1}{2},\quad
c = 1-\frac{\sqrt{3}}{2},
\label{eq:extrap-q4}
\end{equation}

con $a \approx 1{,}866$ y $c \approx 0{,}134$: cada nodo recibe el peso $a$ del punto de Gauss más próximo, $c$ del diagonalmente opuesto y $b$ de los otros dos. La numeración no es un detalle menor: permutar los puntos de Gauss permuta las columnas de $\bm{E}$, y una matriz escrita para otra numeración corrompe las tensiones nodales sin alterar los desplazamientos, de modo que una verificación basada solo en desplazamientos no lo detecta. Por eso el motor construye $\bm{E}$ invirtiendo numéricamente $\bm{M}$ con sus propios puntos de Gauss, para Q4 igual que para Q9, y una prueba de regresión comprueba que $\bm{E}$ reproduce exactamente cualquier campo bilineal (o bicuadrático) muestreado en esos puntos. Para el Q9, $\bm{M}$ es de $9\times9$ y su inversa no admite forma cerrada simple. La extrapolación se aplica a las componentes cartesianas $(\sigma_x,\sigma_y,\tau_{xy})$; las tensiones principales y la de von Mises, al ser funciones no lineales de ellas, se recalculan a partir de las componentes ya extrapoladas en lugar de extrapolarse de forma directa: es el criterio que adopta EduFEM para que cada valor nodal sea consistente con sus propias componentes cartesianas.

Finalmente, como un nodo compartido recibe valores distintos desde cada elemento adyacente, se aplica un \emph{promediado nodal}

\begin{equation}
\bm{\sigma}_{\text{nodo}} = \frac{1}{n_e}\sum_{e=1}^{n_e} \bm{\sigma}_{\text{nodo}}^{(e)},
\label{eq:promediado}
\end{equation}

que suaviza el campo y produce un contorno continuo. El salto entre los valores sin promediar es, además, un indicador útil del error de discretización —es el principio en que se apoyan los estimadores de error por recuperación \autocite[pp.~476-477]{zienkiewicz2013fem}—: cuanto mayor la dispersión entre elementos, más gruesa es la malla en esa zona.

\section{Calidad de malla}
\label{sec:calidad-malla}

La precisión del resultado depende de la calidad geométrica de los elementos. EduFEM evalúa dos métricas ortogonales tomadas de la biblioteca de calidad geométrica Verdict \autocite{stimpson2007verdict}. El \textbf{Jacobiano escalado} mide la distorsión angular normalizando el determinante del Jacobiano por las normas de sus filas —los vectores tangentes a las líneas $\xi$ y $\eta$—,

\begin{equation}
q_{SJ} = \min_g \frac{\det\bm{J}_g}{\|\bm{J}_{1,:}\|_g\,\|\bm{J}_{2,:}\|_g},
\label{eq:scaled-jacobian}
\end{equation}

cantidad que equivale al seno del ángulo local entre ambas tangentes y vive en el rango $[-1, 1]$: vale $1$ en un mapeo ortogonal, indica baja calidad cerca de $0$ y delata inversión cuando $q_{SJ} \le 0$. El validador de salud y la memoria de cálculo la evalúan en los puntos de Gauss del elemento; el módulo M0 emplea la forma por vértices de Verdict —el seno del ángulo interior en cada esquina, tomando el peor—, que coincide con la fórmula que el alumno deriva a mano. La \textbf{compacidad} o \emph{stretch} mide la proporción del elemento relacionando su arista más corta $L_{\min}$ con su diagonal más larga $D_{\max}$,

\begin{equation}
Q = \sqrt{2}\,\frac{L_{\min}}{D_{\max}} \in [0, 1],
\label{eq:stretch}
\end{equation}

que vale $1$ para un cuadrado perfecto y disminuye al alargarse o aplastarse el elemento. De la documentación de Verdict provienen los mínimos admisibles del cuadrilátero, $q_{SJ} \ge 0{,}30$ y $Q \ge 0{,}25$. EduFEM eleva el corte del Jacobiano escalado a $q_{SJ} \ge 0{,}50$ —el valor que Verdict fija para el triángulo y el hexaedro— por prudencia pedagógica: en un curso introductorio conviene señalar como inadmisible un elemento que la biblioteca todavía toleraría. El corte superior con que el módulo M0 distingue los elementos «buenos», $q_{SJ} \ge 0{,}80$ y $Q \ge 0{,}50$, es asimismo una convención de la interfaz de EduFEM. El validador de salud y la memoria de cálculo asignan además a cada elemento un estado global (buena, aceptable o mala) combinando cuatro métricas. Las dos primeras salen del propio Jacobiano: $q_{SJ}$ evaluado en los puntos de Gauss y la razón $R_J = \min_g \det\bm{J}_g / \max_g \det\bm{J}_g$ entre el menor y el mayor de sus determinantes. Las dos restantes se construyen sobre una descomposición bilineal del cuadrilátero, equivalente a la que Verdict emplea para sus métricas de forma \autocite{stimpson2007verdict}, que reduce las coordenadas de las cuatro esquinas $\bm{r}_1,\dots,\bm{r}_4$ a tres vectores característicos:

\begin{equation}
\begin{aligned}
\bm{X}_1 &= \tfrac{1}{4}\left(-\bm{r}_1 + \bm{r}_2 + \bm{r}_3 - \bm{r}_4\right), \qquad
\bm{X}_2 = \tfrac{1}{4}\left(-\bm{r}_1 - \bm{r}_2 + \bm{r}_3 + \bm{r}_4\right), \\[2pt]
\bm{X}_3 &= \tfrac{1}{4}\left(\bm{r}_1 - \bm{r}_2 + \bm{r}_3 - \bm{r}_4\right), \\[4pt]
AR &= \frac{\max(\|\bm{X}_1\|, \|\bm{X}_2\|)}{\min(\|\bm{X}_1\|, \|\bm{X}_2\|)}, \qquad
T_R = \max\left(\frac{|\bm{X}_3 \cdot \bm{X}_1|}{\|\bm{X}_1\|^2},\; \frac{|\bm{X}_3 \cdot \bm{X}_2|}{\|\bm{X}_2\|^2}\right).
\end{aligned}
\label{eq:descomposicion-bilineal}
\end{equation}

Los dos primeros vectores unen los puntos medios de las aristas opuestas en cada dirección natural y miden la extensión del elemento; el tercero recoge su parte propiamente bilineal. De ellos se derivan la \textbf{relación de aspecto} $AR$, que vale $1$ en un cuadrado o un rombo y crece al alargarse el elemento, y el \textbf{estrechamiento} (\emph{taper}) $T_R$, que es nulo en todo paralelogramo y crece con la trapecialidad: precisamente el modo de distorsión ante el que el Q4 es más sensible, como muestra el caso de Cook del \autoref{cap:resultados}. En el Q9 la cuarta métrica no es el estrechamiento, sino la desviación de los nodos de arista y central respecto de su posición ideal sobre el elemento macro. Los cortes con que se asigna el estado global son convención de EduFEM y no umbrales bibliográficos: se exige $q_{SJ} \ge 0{,}7$, $R_J \ge 0{,}5$, $AR \le 3$ y $T_R \le 0{,}3$ —o desviación $\le 0{,}10$ en el Q9— para «buena», y $0{,}3$, $0{,}2$, $5$ y $0{,}5$ —o $0{,}25$ en el Q9— para «aceptable». Las dos métricas de M0 son complementarias: un rectángulo muy alargado puede tener $q_{SJ} = 1$ (ángulos rectos) y, aun así, una compacidad inadmisible que degrada la precisión.

\section{Verificación y validación}
\label{sec:verificacion-validacion}

Un programa de cálculo debe distinguir entre \emph{verificación} —comprobar que las ecuaciones se resuelven correctamente— y \emph{validación} —comprobar que las ecuaciones representan la realidad física \autocite[pp.~13-14]{oberkampf2010vv}. EduFEM incorpora ambas estrategias. Conviene precisar el alcance del segundo término: en su acepción estricta la validación exige el contraste con mediciones experimentales, y este trabajo no realiza campaña experimental propia. Sus referencias externas son una solución analítica de la teoría de la elasticidad y un modelo equivalente resuelto en un programa comercial, de modo que el contraste corresponde a verificación de solución y a comparación código a código; esta última, según la propia literatura, no sustituye a la verificación rigurosa del código \autocite{oberkampf2010vv}.

La verificación de código se realiza con el \textbf{Método de las Soluciones Manufacturadas} (MMS): se elige a priori un campo de desplazamientos exacto $\bm{u}_M$, se sustituye en las ecuaciones de equilibrio para deducir la fuerza de cuerpo $\bm{b} = -\nabla\cdot\bm{\sigma}(\bm{u}_M)$ que lo produce, y se inyectan esa carga y las restricciones de Dirichlet no homogéneas correspondientes \autocites[pp.~21 y 26]{salari2000mms}[pp.~219-222]{oberkampf2010vv}. El error de la solución numérica frente a la exacta se cuantifica con las normas $L^2$ del desplazamiento y $H^1$ de su gradiente:

\begin{equation}
\begin{aligned}
\|\bm{u}_h - \bm{u}\|_{L^2} &= \left( \int_\Omega \|\bm{u}_h - \bm{u}\|^2 \, d\Omega \right)^{1/2}, \\
|\bm{u}_h - \bm{u}|_{H^1} &= \left( \int_\Omega \|\nabla\bm{u}_h - \nabla\bm{u}\|^2 \, d\Omega \right)^{1/2}.
\end{aligned}
\label{eq:normas-error}
\end{equation}

Al refinar la malla, estas normas deben decrecer con las tasas asintóticas teóricas ($\mathcal{O}(h^{p+1})$ en $L^2$ y $\mathcal{O}(h^p)$ en $H^1$, con $p$ el orden polinómico del elemento) que establece la teoría de aproximación del MEF \autocites[pp.~106-107]{strang2008analysis}[pp.~189-190]{hughes2000fem}, lo que confirma la correcta implementación. Las integrales de error se evalúan con un punto de Gauss más por dirección que la rigidez ($3\times3$ para Q4, $4\times4$ para Q9) para no subestimar artificialmente el error: es una precaución de implementación de los propios guiones de verificación, porque el error de la cuadratura con que se integra la norma puede interferir con el error de discretización y degradar el orden observado. Como esas normas dependen solo del campo de desplazamientos, no ejercitan la cadena de recuperación de tensiones (evaluación en los puntos de Gauss, extrapolación a los nodos y promediado); por ello se mide además la norma $L^2$ del campo de tensiones recuperado $\bm{\sigma}^*_h$ frente a $\bm{\sigma}(\bm{u}_M)=\bm{D}\,\bm{\varepsilon}(\bm{u}_M)$ —el error de las magnitudes derivadas debe verificarse por separado, y su orden puede ser menor que el de la solución primaria \autocite[p.~29]{salari2000mms}—, y se comprueba de forma independiente que la matriz de extrapolación reproduce exactamente los campos polinómicos de su espacio de interpolación.

La validación contrasta el motor contra problemas de referencia con solución conocida. La \textbf{viga de Timoshenko} simplemente apoyada admite solución analítica de la teoría de la elasticidad \autocite[art.~21, pp.~39-43]{timoshenko1970elasticity}, y EduFEM la reproduce con errores inferiores al $0{,}3\,\%$ en las magnitudes primarias —tensión normal y flecha—, que el \autoref{cap:resultados} cuantifica en detalle, comparando además contra un modelo de cáscara de SAP2000 \autocite{csi2017sap2000}. La \textbf{membrana de Cook}, un trapecio en voladizo con carga de corte que debe su nombre a R. D. Cook, es el banco de prueba clásico de la literatura para evaluar la sensibilidad de los elementos a la distorsión geométrica y al bloqueo por cortante: el desplazamiento del extremo converge al valor de referencia de uso convencional, $23{,}96$, pero el elemento Q4 lo subestima sensiblemente en las mallas gruesas por el bloqueo por cortante (\emph{shear-locking}) descrito en la \autoref{sec:locking}; el elemento Q9, de orden superior, lo evita y converge con rapidez \autocites[p.~98]{cook2002concepts}[p.~221]{hughes2000fem}. Estos contrastes muestran por qué la elección del tipo de elemento condiciona la calidad del resultado.
